% ===========================================================================
% Chapter 10: Conclusion
% ===========================================================================
\chapter{Conclusion}
\label{ch:conclusion}

\section{Key Findings}
\label{sec:key-findings}

\textbf{The model's absorber had been switched off, and switching it on locates its
effect.} The absorber term of VULCAN-Venus was present but inactive. Once active, it
changes the radiation field in proportion to its amount, and the change lives in one band
of altitude, 69 to 93~km.

\textbf{Where an absorber sits matters more than its spectrum.} The same absorber has 7.34
times more effect moved up than moved down, against a factor of 4.95 for its spectrum at
a fixed height.

\textbf{The chemistry decides where a candidate sits.} Placed by the model's own
chemistry, ClS\sub{2}, the gas closest in amount, falls below the response band and
becomes radiatively the same as no absorber. Being made in the right amount and absorbing
at the right wavelengths is not enough; a candidate must also be in the right place.

\textbf{No species in the model passes all three tests.} Of 23 species that could be
screened, none supplies enough absorption, and the three tests disagree about which comes
closest.

\textbf{Two published candidates are reassessed.} The proposed iron and sulfur minerals
cannot absorb enough at 375~nm even as pure solids. The published rejection of ferric
chloride rests on a measurement in water; if ferric chloride absorbs in the droplets as in
hydrochloric acid, it is a plausible minor component, though its altitude is untested.

\textbf{Every number is verified.} The registry recomputes every quantity from the saved
output, the checker tests every printed number precise enough to test, and
Chapter~\ref{ch:verify} traces numbers to the code.

Phase~2 asks whether these verdicts survive the much stronger mixing now estimated for
the Venus clouds, and whether it can carry a candidate into the band.
